%% =====================================================================
%%  T-22-01  The Liking Handle
%%  -------------------------------------------------------------------
%%  One idea per file. Core is mandatory. Slots in canonical order.
%%  Max 700 words. No layout commands. No filler. New source material
%%  for this entry goes in sources/<BOOK>/T-22-01.tex -- never by
%%  rewriting this file (docs/RULES.md R0.1).
%% =====================================================================
%% @kind: topic
%% @id: T-22-01
%% @title: The Liking Handle
%% @subtitle: Agreement is bought more cheaply by being liked than by being right
%% @chapter: c22
%% @order: 10
%% @status: stable
%% @sources: B4
%% @updated: 2026-09-19
%% @allow-rewrite: slot migration to the seven-question format (docs/RULES.md section 2)

\topic{T-22-01}{The Liking Handle}{Agreement is bought more cheaply by being liked than by being right}


\begin{core}
People say yes to those they like, and the effect is large enough that liking can be manufactured more cheaply than a case can be argued. Five things produce it: resemblance, praise, familiarity built through working together, surface appeal, and association with good news.
\end{core}

\begin{mechanism}
Liking is used as a shortcut for trustworthiness, on the reasonable assumption that someone who is like you, who approves of you, and whom you have cooperated with is unlikely to harm you. The assumption is sound about people you know. It is not sound about people who have studied the cues.

Each component can be produced independently of the quality it stands for. Resemblance is a matter of matching posture, vocabulary, pace and background, all of which are observable and copyable. Praise requires only that it be pleasant to receive --- and it remains pleasant when the recipient knows it is strategic. Familiarity is manufactured by arranging repeated small cooperation rather than by waiting for a relationship to develop. Surface appeal is dressing, grooming and setting. Association is attaching yourself to whatever the counterpart already likes. None of the five requires sincerity, and none is switched off by being recognised.
\end{mechanism}

\begin{application}
  \item Establish resemblance before making the request: pace, vocabulary, and one genuine point of common background.
  \item Praise specifically and early. Vague praise reads as technique; specific praise reads as attention.
  \item Arrange a small cooperative task before the real negotiation, so that the relationship exists when the disagreement arrives.
  \item Manage the setting and your own presentation; surface appeal is cheap and is weighted far above its information content.
  \item Attach the proposal to something the counterpart already values rather than presenting it on its own.
\end{application}

\begin{feedback}
  \item You find the person agreeable and cannot say what they have actually agreed to.
  \item They resemble you in ways that became visible only after you met.
  \item They praise you early, specifically, and before asking anything.
  \item A small joint task is proposed before the real one.
  \item You are more willing to concede to them than the merits warrant.
\end{feedback}

\begin{failure}
The components are visible to anyone who knows them, and a person caught manufacturing resemblance loses the credibility of everything genuine they had. It also degrades the practitioner: someone who  habitually matches, praises and attaches themselves to good news eventually finds it hard to tell which of their own opinions are theirs, which is a real loss and is described in Part VII.
\end{failure}

\begin{countermeasures}
  \item Separate the person from the proposal. Ask whether you would accept the same terms from someone you found disagreeable.
  \item Discount resemblance that appeared after the meeting began. Matching is a technique; shared history is not.
  \item Receive praise and do not repay it. Being approved of creates an obligation to approve back, which is the point of it.
  \item Decline the small joint task, or complete it and then re-open the real question from scratch.
  \item Delay. Liking decays faster than a good argument, and the decision improves overnight.
\end{countermeasures}

\sources{B4}
\seesources
\seealso{T-22-02, T-23-02, T-04-03, T-24-02}
